\documentclass[12pt]{article}
\usepackage{amsfonts,amsthm,amsmath,amssymb,mathtools}
\usepackage{mathrsfs}
\usepackage{enumitem}
\usepackage{tikz}
\usepackage{tikz-cd}
\usepackage{geometry}
\usepackage{mlmodern}

\usetikzlibrary{shapes}

\geometry{letterpaper, portrait, margin=1in}

\setcounter{section}{5}

\DeclareMathOperator{\range}{range}
\DeclareMathOperator{\sgn}{sgn}

% differential operator
\newcommand{\dd}{\mathop{}\!\mathrm{d}}

\newtheorem{thm}{Theorem}
\newenvironment{theorem}[1]{
	\renewcommand{\thethm}{#1}
	\begin{thm}
		}{\end{thm}}

\newtheorem{lma}{Lemma}
\newenvironment{lemma}[1]{
	\renewcommand{\thelma}{#1}
	\begin{lma}
		}{\end{lma}}

\theoremstyle{definition}
\newtheorem{ex}{}[section]
\newenvironment{exercise}[1]{
  \renewcommand{\theex}{\arabic{section}.\arabic{ex}}
	\begin{ex}
		}{\end{ex}}

\title{\S\arabic{section} Exercises}
\author{Connor Mooneyhan}
\date{May 31, 2025}

\begin{document}

\maketitle

\begin{ex}
	\textbf{The real line with two origins} \\
	Let $A$ and $B$ be two points not on the real line $\mathbb{R}$.
	Consider the set $S = (\mathbb{R} - \left\lbrace 0 \right\rbrace) \cup \left\lbrace A, B \right\rbrace$.
  \begin{center}
    \begin{tikzpicture}
      \draw[thick] (-2, 0) -- (-2.5pt, 0);
      \draw[thick] (2.5pt, 0) -- (2, 0);
      \draw[fill=black] (0, .4) circle (2.5pt) node [anchor=south west] {A};
      \draw[fill=black] (0, -.4) circle (2.5pt) node [anchor=north west] {B};
    \end{tikzpicture}
  \end{center}
\end{ex}

\end{document}
